% ====================================================================
% FILE: content/14_disciplines_extended.tex
% Кросс-дисциплинарные инстанцирования континуумной рамки
% Core 1.3.3 flagship version
% ====================================================================

\section{Кросс-дисциплинарные инстанцирования континуумной рамки}
\label{sec:disciplines-extended}

В этой главе проверяется ограничительное утверждение о представимости:
выбранные научные и формально-теоретические области могут быть выражены
кортежем
\[
  (\Omega, A, P(t), J(t), \Theta, \partial\Omega, C, k(t))
\]
таким образом, что это совместимо со структурной иерархией уровней
\(K_0\)–\(K_{10}\), ландшафтом порогов, семейством операторов и топологией
ветвей раздела~the branching-topology chapter.
Здесь \(k(t)\) не содержит индекса уровня только для удобства чтения кортежа;
локально это означает соответствующее значение \(k(K_x,t)\).
Аналогично, в последующих блоках предметной области может опускаться
параметр \(t\) в обозначениях вроде \(P_x\) или \(J_x\), когда описывается
семейство потенциалов или потоков, а не срез в фиксированный момент времени.
Каждая предметная карта формулируется здесь под единым каноническим ярлыком:
\emph{правилом кандидирующего вложения}. Сочетания «тест ограниченной
представимости», «критерий структурной совместимости», «проход ограниченного
вложения» и «критерий ограниченного вложения» являются локальными
формулировками именно этого правила, а не самостоятельными независимыми
критериями. Цель — показать, насколько хорошо описанные классы систем
представимы как континуумы в рамках структурных аксиом OC, а не утверждать,
что каждый объект дисциплины исчерпывающим образом описывается одной
энкодингом OC \parencite{anderson_more_is_different}.

Это структурное требование совместимости, а не утверждение об
операционной выгоде. Практическое сравнение отложено на
раздел~the oc core 1 3 practical discussionutilityи
приложение~the practical utility chapter atlas.
Эта глава спрашивает, может ли область быть представлена как ограниченный
континуум; она не утверждает, что у каждой такой области уже есть замкнутый
пакет теорем, отложенный аудит, сравнительный тест или готовая к внедрению
процедура.

% ====================================================================
\subsection{Принципы вложения предметной области}
% ====================================================================

Чтобы представить дисциплину внутри OC, должны выполняться следующие
условия структурной согласованности.

\paragraph{1. Отображение объектов на компоненты континуума.}
Дисциплина имеет кандидатное континуумное представление, если её объекты
можно сопоставить:
\begin{itemize}
  \item \(\Omega\): допустимые конфигурации,
  \item \(A\): оси, задающие существенные степени свободы,
  \item \(P\): потенциалы, кодирующие силы, энергии, стоимости или
        семантическое напряжение,
  \item \(J\): потоки (вещество, заряд, вероятность, информация,
        ресурсы),
  \item \(\Theta\): пороги, управляющие устойчивостью, фазовым переходом и
        схлопыванием,
  \item \(\partial\Omega\): границы жизнеспособности,
  \item \(C\): циклы, поддерживающие персистентность,
  \item \(k(t)\): континуумичность (мера жизнеспособности).
\end{itemize}

\paragraph{2. Ограничения пространства вложения.}
Допустимое инстанцирование требует одновременного выполнения
\(\Omega(K_x)\subseteq\Omega(M_x)\) и
\(A(K_x)\subseteq A(M_x)\). Если требуемые состояния или оси отсутствуют
в пространстве вложения, предлагаемое инстанцирование структурно недопустимо.

\paragraph{3. Соответствие уровню K.}
Дисциплина ограниченно вложена на уровне \(K_x\) только тогда, когда:
\begin{itemize}
  \item доминирующие оси соответствуют осям \(K_x\);
  \item пороги уточняются относительно более низких уровней;
  \item потоки и циклы согласованы с семейством операторов \(K_x\);
  \item поведение удовлетворяет размерностной монотонности.
\end{itemize}

\paragraph{4. Структурная совместимость.}
Инстанцирование области требует:
\begin{itemize}
  \item внутренней согласованности с порогами OC,
  \item существования по крайней мере одного поддерживающего цикла
        \(C_{\mathrm{support}}\),
  \item непустой допустимой области \(\Omega\neq\emptyset\),
  \item расширяемости динамики области через семейство операторов.
\end{itemize}

Кандидатное вложение отвергается при любом срыве проверки: допустимая
область пустая, требуемые оси отсутствуют в пространстве вложения, заявленные
пороги не могут быть одновременно выполнены или предлагаемые циклы не удерживают
объект от пересечения соответствующей границы. Это критерии отказа, а не
просто косметические пробелы.

Следующие подразделы суммируют ограничительные кандидатные представления
для разных дисциплин.

% ====================================================================
\subsection{Физика (уровень \texorpdfstring{$K_2$}{K2})}
% ====================================================================

Физические системы реализуют уровень \(K_2\), когда конфигурации
поля, параметры порядка и структуры когерентности задают соответствующие оси.
Обсуждение намеренно ограничено стандартными режимами критических
явлений, ренормгруппового подхода и перколяции, а не безграничным
утверждением о всей физике \parencite{goldenfeld_lectures,wilson_renormalization,stauffer_percolation}.

\paragraph{Пространство состояний и оси.}
\begin{itemize}
  \item \(\Omega_{2}\): конфигурации полей (скалярные, векторные,
        калибровочные поля).
  \item \(A_{2}\): пространственные оси, модовые оси, оси внутренних симметрий,
        оси параметров порядка.
\end{itemize}

\paragraph{Потенциалы и потоки.}
\begin{itemize}
  \item \(P_2\): физические функционалы энергии (действие, свободная энергия).
  \item \(J_2\): диффузионные, конвективные и полевыя токи.
\end{itemize}

\paragraph{Пороги.}
\begin{itemize}
  \item \(\Theta_{\mathrm{crit}}\): границы фаз,
  \item \(\Theta_{\mathrm{BKT}}\): порог развязывания вихрей
        Бере́зинского\textendash\mbox{Костерлица\textendash}Тоу,
  \item \(\Theta_{\mathrm{mass}}\): порог когерентности для возникновения массы,
  \item \(\Theta_{\mathrm{death}}\): условия конфайнмента или декуплинга.
\end{itemize}

\paragraph{Граничные структуры.}
В этой ограничительной постановке \(\partial\Omega\) соответствует критическим
поверхностям, границам перколяции и интерфейсам фазового разделения.

\paragraph{Циклы.}
Топологические солитоны, пары вихрей, ренормгрупповые циклы и
осциллирующие моды поля являются примерами \(C_2\).

\paragraph{Критерий ограничительного вложения (Физика → \texorpdfstring{$K_2$}{K2}).}
Физическая система удовлетворяет критерию вложения \(K_2\), если:
\begin{itemize}
  \item её поля определяют связное конфигурационное пространство \(\Omega_2\),
  \item функционал энергии задаёт пороги, совместимые с OC,
  \item потоки подчиняются законам сохранения, выражаемым через \(J_2\),
  \item когерентные циклы поддерживают структурную персистентность.
\end{itemize}

% ====================================================================
\subsection{Химия и происхождение жизни (уровни
            \texorpdfstring{$K_3$}{K3}--\texorpdfstring{$K_4$}{K4})}
% ====================================================================

Химические системы описываются на уровне \(K_3\), а протоклеточные — на
уровне \(K_4\).
Эти отображения ограничены замыканием реакционных сетей, каталитической
организацией и мембранно-структурной пребиотической ограниченностью
\parencite{kauffman_autocatalytic,raf_networks}.

\paragraph{Химические континуумы (\texorpdfstring{$K_3$}{K3}).}
\begin{itemize}
  \item \(\Omega_{\mathrm{chem}}\): векторы концентраций, состояния реакций,
        каталитические конфигурации.
  \item \(A_{\mathrm{chem}}\): концентрации видов, внешние параметры
        (pH, солёность, температура).
  \item \(P_{\mathrm{chem}}(t)\): химические потенциалы, градиенты сродства и
        дисбалансы концентраций.
  \item \(J_{\mathrm{chem}}\): скорости реакций, транспортные потоки.
  \item \(\Theta_{\mathrm{act}}\): порог каталитической активации.
  \item \(\Theta_{\mathrm{closure}}\): порог замыкания рефлексивно
        автокаталитических и пищевых (RAF) сетей.
  \item \(C_{\mathrm{chem}}\): реакционные циклы, циклы замыкания RAF.
\end{itemize}

\paragraph{Протоклеточные континуумы (\texorpdfstring{$K_4$}{K4}).}
\begin{itemize}
  \item \(\Omega_{4}\): состояния пузырьков, конфигурации мембран,
        внутренние химические составы.
  \item \(A_4\): кривизна, поверхностное напряжение, проницаемость,
        заряд.
  \item \(P_4(t)\): осмопотенциалы, химические, электрохимические и
        кривизностные потенциалы.
  \item \(\Theta_{\mathrm{perm}}, \Theta_{\mathrm{osm}},
        \Theta_{\mathrm{curv}}\): пороги жизнеспособности мембран.
  \item \(J_4\): осмотические потоки, ионные потоки, метаболические потоки.
  \item \(C_4\): метаболические циклы, циклы поддержания мембраны,
        циклы устойчивости градиентов.
\end{itemize}

\paragraph{Критерий ограничительного вложения (химия → \texorpdfstring{$K_3$}{K3}).}
Химическая система удовлетворяет критерию вложения \(K_3\), если:
\begin{itemize}
  \item реакционные сети формируют непустую допустимую область,
  \item каталитические и пороги замыкания представимы без мембранных
        специфических осей,
  \item циклы поддерживают замыкание реакций и контроль концентраций,
  \item ограничения вложения поддерживают химические оси, требуемые
        маршрутом.
\end{itemize}

\paragraph{Критерий ограничительного вложения (протоклетки →
            \texorpdfstring{$K_4$}{K4}).}
Протоклеточная система удовлетворяет критерию вложения \(K_4\), если:
\begin{itemize}
  \item мембранно-поддерживаемая химическая динамика формирует
        непустую допустимую область,
  \item границы мембран определяют \(\partial\Omega_4\),
  \item циклы поддерживают концентрационные градиенты и целостность
        компартмента,
  \item ограничения вложения поддерживают как химические, так и мембранные
        оси.
\end{itemize}

% ====================================================================
\subsection{Биология (уровень \texorpdfstring{$K_5$}{K5})}
% ====================================================================

Биологические клетки, сверх протоклеточного \(K_4\)-маршрута, возникают,
когда протоклетки приобретают возбудимость и структурированную регуляцию.
Цель здесь — представление ограниченного режима возбудимых и морфогенетических
организаций, а не неограниченная теория жизни как таковой
\parencite{turing_morphogenesis,hopfield_nn}.

\paragraph{Пространство состояний и оси.}
\begin{itemize}
  \item \(\Omega_5\): распределения потенциала мембраны, состояний каналов,
        ионных градиентов.
  \item \(A_5\): оси возбудимости (потенциал мембраны \(\Delta V\),
        концентрации ионов, оси переменных затвора).
\end{itemize}

\paragraph{Потенциалы и потоки.}
\[
  P_5:\text{ электрохимические потенциалы},\quad
  J_5:\text{ ионные токи, метаболические потоки}.
\]

\paragraph{Пороги.}
\[
  \Theta_{\mathrm{exc}},\quad
  \Theta_{\mathrm{refr}},\quad
  \Theta_{\mathrm{grad}},\quad
  \Theta_{\mathrm{death}}.
\]

\paragraph{Циклы.}
Релевантные циклы включают протоспайк-цикл \(C_{\mathrm{spike}}\),
метаболические циклы и петли электрических осцилляций.

\paragraph{Критерий ограничительного вложения (биология → \texorpdfstring{$K_5$}{K5}).}
Биологическая система удовлетворяет критерию вложения \(K_5\), если:
\begin{itemize}
  \item возбудимость задает нетривиальные оси \(\Delta V\),
  \item пороги определяют структуру импульса и рефрактерного периода,
  \item циклы поддерживают гомеостаз и сигнализацию,
  \item вложение включает ионные, энергетические и мембранно-поддерживающие
        оси.
\end{itemize}

% ====================================================================
\subsection{Когниция (уровень \texorpdfstring{$K_6$}{K6})}
% ====================================================================

Ниже уровень \(K_7\) разделы \(K_6\)\textendash\(K_{10}\) дают
структурные кандидатные вложения.
Это эскизы на границе исследовательских направлений, пока
раздел~the oc core 1 3 practical discussionutilityили приложение
the practical utility chapter atlasне выделит более узкий
операционный канал.
Самостоятельно они не являются практическими утверждениями для непосредственного
использования.

Когнитивные системы возникают, когда представления, связывание и
предиктивные модели становятся доминирующими осями.
Обсуждение опирается на классические работы по организации поведения и
мотивам предиктивной обработки, а не на тезис о том, что OC полностью
заменяет когнитивную науку \parencite{hebb_organization,friston_free_energy}.

\paragraph{Пространство состояний и оси.}
\begin{itemize}
  \item \(\Omega_6\): конфигурации представлений, активированные
        концептуальные состояния, состояния предсказания.
  \item \(A_6\): оси представления, оси связывания, предиктивные оси.
\end{itemize}

\paragraph{Потенциалы и потоки.}
\begin{itemize}
  \item \(P_6\): ошибка предсказания, семантическое напряжение,
        уверенность.
  \item \(J_6\): потоки внимания, потоки вывода, переходы памяти.
\end{itemize}

\paragraph{Пороги.}
\[
  \Theta_{\mathrm{pred}},\quad
  \Theta_{\mathrm{bind}},\quad
  \Theta_{\mathrm{coh}},\quad
  \Theta_{\mathrm{expr}}.
\]

\paragraph{Циклы.}
Релевантные циклы включают циклы внимания\textendash предсказания,
циклы консолидации памяти и циклы устойчивости модели.

\paragraph{Критерий ограничительного вложения (когниция →
            \texorpdfstring{$K_6$}{K6}).}
Когнитивная система удовлетворяет критерию вложения \(K_6\), если:
\begin{itemize}
  \item она поддерживает когерентные внутренние модели в рамках порогов
        выражаемости,
  \item потоки снижают ошибку предсказания и поддерживают концептуальную
        структуру,
  \item циклы регулируют внимание, память и обновление,
  \item границы \(\partial\Omega_6\) кодируют ограничения модели.
\end{itemize}

% ====================================================================
\subsection{Общество (уровень \texorpdfstring{$K_7$}{K7})}
% ====================================================================

Социальные системы реализуют \(K_7\), когда доминируют нормы, институты
и доверие как структурные оси.
Предметная область — пороговое коллективное поведение, сетевой
координационный режим и институциональная персистентность
\parencite{granovetter_threshold,schelling_micromotives,newman_networks}.

\paragraph{Пространство состояний и оси.}
\begin{itemize}
  \item \(\Omega_7\): коммуникационные графы, ролевые структуры,
        институциональные состояния.
  \item \(A_7\): нормативные оси, оси доверия, институциональные оси.
\end{itemize}

\paragraph{Потенциалы и потоки.}
\[
\begin{aligned}
  P_7 &: \text{доверие, легитимность, символический капитал},\\
  J_7 &: \text{коммуникация, обмен ресурсами, передача норм}.
\end{aligned}
\]

\paragraph{Пороги.}
\[
  \Theta_{\mathrm{trust}},\quad
  \Theta_{\mathrm{leg}},\quad
  \Theta_{\mathrm{role}}.
\]

\paragraph{Циклы.}
Релевантные циклы включают институциональные циклы, нормативные циклы и
циклы координации.

\paragraph{Критерий ограничительного вложения (общество →
            \texorpdfstring{$K_7$}{K7}).}
Социальная система удовлетворяет критерию вложения \(K_7\), если:
\begin{itemize}
  \item доверие и нормы задают пороги жизнеспособности,
  \item институты поддерживают циклы управления и координации,
  \item потоки коммуникации и ресурсов соблюдают ограничения вложения,
  \item социальные границы формируют непустую \(\partial\Omega_7\).
\end{itemize}

% ====================================================================
\subsection{Цивилизационные системы (уровень
            \texorpdfstring{$K_8$}{K8})}
% ====================================================================

Цивилизационные системы расширяют социальные континуумы за счёт
инфраструктур, энергетических потоков и масштабной координации.
Область ограничивается крупномасштабными организованными системами с
сетевой инфраструктурой и ограничениями устойчивости
\parencite{newman_networks,may_stability_complexity}.

\paragraph{Пространство состояний и оси.}
\begin{itemize}
  \item \(\Omega_8\): инфраструктурные сети, экономические состояния,
        распределения населения.
  \item \(A_8\): инфраструктурные оси, оси потоков энергии,
        оси коммуникации, оси сложности.
\end{itemize}

\paragraph{Потенциалы и потоки.}
\[
  P_8: \text{доступность энергии, ресурсные градиенты, потенциальные риски},\quad
  J_8: \text{торговля, энергетические потоки, потоки данных}.
\]

\paragraph{Пороги.}
\[
  \Theta_{\mathrm{stress}},\quad
  \Theta_{\mathrm{capacity}},\quad
  \Theta_{\mathrm{cohesion}}.
\]

\paragraph{Циклы.}
Релевантные циклы включают экономические циклы, циклы обновления
инфраструктуры и петли трансфера знания.

\paragraph{Критерий ограничительного вложения (цивилизационные системы →
            \texorpdfstring{$K_8$}{K8}).}
Цивилизационная система удовлетворяет критерию вложения \(K_8\), если:
\begin{itemize}
  \item энергетические и ресурсные структуры задают \(P_8\),
  \item инфраструктурные потоки формируют устойчивые глобальные циклы,
  \item институты уровня \(K_7\) последовательно встраиваются как
        подподузлы,
  \item крупномасштабные границы определяют глобальную
        \(\partial\Omega_8\).
\end{itemize}

% ====================================================================
\subsection{Теория и формальные системы (уровни
            \texorpdfstring{$K_9$}{K9}--\texorpdfstring{$K_{10}$}{K10})}
% ====================================================================

Теоретический и метатеоретический уровни занимают верхний маршрут, рассматриваемый
в этой главе. Этот подраздел даёт OC-внутреннюю схему для
\(K_9\)\textendash\(K_{10}\), с локальной опорой разделами
the klevels discussionfullи~the foundational-consistency chapter.
Не утверждается о полностью замкнутой внешней теореме вложения для каждой
теории или формальной системы. Рекурсивные структурные континуумы остаются на
уровне \(K_{11}\), вне текущего правила инстанцирования.

\paragraph{Теоретические континуумы (\texorpdfstring{$K_9$}{K9}).}
\begin{itemize}
  \item \(\Omega_9\): пространство моделей, формализмов и представлений.
  \item \(A_9\): концептуальные, формальные и моделирующие оси.
  \item \(P_9\): потенциалы когерентности, эмпирического напряжения,
        объяснительной силы.
  \item \(J_9\): потоки пересмотра моделей, обработки аномалий и
        переноса аргументации.
  \item \(\Theta_{\mathrm{expr}}, \Theta_{\mathrm{coh}}\): пороги
        согласованности и выражаемости.
  \item \(\partial\Omega_9\): границы несогласованности, провала
        выражаемости или неремонтированной аномалии.
  \item \(C_9\): парадигматические циклы, исследовательские циклы.
  \item \(k_9(t)\): континуумичность полосы согласования и валидации.
\end{itemize}

\paragraph{Метатеоретические формальные системы (\texorpdfstring{$K_{10}$}{K10}).}
\begin{itemize}
  \item \(\Omega_{10}\): пространство метатеоретических состояний, формальных
        систем, состояний доказательств, объектов «модель-модели» и
        языков с ограниченной рекурсией.
  \item \(A_{10}\): оси метатеории, рекурсии, согласованности,
        выражаемости и сравнения теорий.
  \item \(P_{10}\): потенциал емкости доказательных потоков,
        запаса согласованности, формальной выразительной емкости и
        потенциальной метатеоретической когерентности.
  \item \(J_{10}\): потоки поиска доказательств, формального обновления,
        переводов метамоделей и потоков управления рекурсией.
  \item \(\Theta_{\mathrm{cons}}, \Theta_{\mathrm{rec}},
        \Theta_{\mathrm{meta}}\): пороги согласованности, ограниченной
        рекурсии и метауровневой когерентности.
  \item \(\partial\Omega_{10}\): несогласованность, недоказуемое
        переразрастание, поверхности рекурсивной границы или провалы
        метатеоретической границы.
  \item \(C_{10}\): циклы доказательств, циклы рекурсии,
        циклы формального обновления и циклы аудита «модель-модели».
  \item \(k_{10}(t)\): формальная и метатеоретическая жизнеспособность
        при ограничениях согласованности, рекурсии, самоссылки и пропускной
        способности доказательств.
\end{itemize}

\paragraph{Кандидатное правило вложения (теория →
            \texorpdfstring{$K_9$}{K9}, метатеоретическая формальная система →
            \texorpdfstring{$K_{10}$}{K10}).}
Теоретический и метатеоретический формальные маршруты удовлетворяют
кандидатному правилу вложения, когда соответствующая структура \(K_9\) или
\(K_{10}\) выполняет следующие условия:
\begin{itemize}
  \item их концептуальные и формальные оси образуют валидное множество осей,
  \item ограничения согласованности задают пороги, совместимые с OC,
  \item циклы рассуждения, моделирования и пересмотра поддерживают
        жизнеспособность,
  \item пространства вложения поддерживают требуемые оси выражаемости.
\end{itemize}

% ====================================================================
\subsection*{Итог}

В этой главе формулируются ограничительные кросс-дисциплинарные критерии
представимости по областям физики, химии, протоклеток, биологии,
когниции, общества, цивилизационных систем, теории и формальных систем.
Каждый ограниченный маршрут отождествляется с одним и тем же структурным
кортежем на детальном уровне, заданном в соответствующем подразделе, но
сохраняется утверждение о структурной совместимости, а не универсальной
подмене локальных моделей каждой дисциплины в OC. Более сильный
операционный и сравнительный вопрос рассматривается позднее в явно указанных
поддерживающих классах, перечисленных
в разделе~the oc core 1 3 practical discussionutilityи атласе строк данных
приложения~the practical utility chapter atlas.

% END OF FILE
